\chapter{Literature Review}
\label{ch:lit_rev} 

Given this project's nature as an implementation-focused project rather than a research-based project, our literature review will be mainly focused towards describing the background and context surrounding this project.

\section{Distributed Systems} 
\subsection{Consensus}
Consensus is a fundamental problem in distributed computing, which asks how to enable a group of nodes to agree on a single shared value \cite{dsbook}. A consensus protocol is a mechanism for all nodes to communicate in a known manner such that consensus may be achieved within the network.
\\

Consensus is essential in enabling techniques like State Machine Replication, where a single service is replicated to achieve greater redundancy and throughput \cite{smr}. In this case, consensus is required to ensure that the internal state of the service remains consistent across all its replicas. 
\\

\subsection{Byzantine Fault-Tolerant Consensus}
Consensus can be made harder to achieve under certain scenarios, where communication between nodes can be lost, or worse, manipulated by an adversary. In such scenarios, traditional consensus protocols may no longer suffice to ensure consensus within the distributed system. One such extension to the consensus problem is the Byzantine Fault-Tolerant (BFT) consensus problem, where consensus protocols must now tolerate "Byzantine faults", or situations where nodes may behave in an arbitrary manner, such as by sending out-of-date information or even acting in a malicious way \cite{bgp-lamport}.
\\

BFT consensus protocols are commonly used in the industry, most notably in blockchain-related applications. For example, the HotStuff protocol is used in SafeStake \cite{safestake}, and the Tendermint protocol is used in the Cosmos Network \cite{cosmos}.
~\\

\section{Formal Verification}
\subsection{Introduction}
Formal Verification is a method used to mathematically prove that a given system meets certain specifications or fulfills certain properties. This is done by modelling the system as mathematical structures, allowing one to rigorously prove the required claims \cite{fvdefn}. Formal verification is an important technique to guarantee correctness in mission-critical situations. One example is the CompCert C compiler, which guarantees that machine code generated by the compiler is effectively equivalent to the original source code \cite{compcert}.
\\

However, formal verification of software can be complex and expensive \cite{fv-disc}. An additional problem is ensuring that the mathematical model created is a faithful representation of the original system, otherwise, the proven properties on the modelled system would not necessarily hold on the actual system. Likewise, the veracity of the proofs used are also crucial: if an error exists in the proof of a certain property, we cannot actually guarantee that such a property holds for the original system.
\\

Another potential limitation of formal verification techniques is that the proven properties may only hold for a particular version of the specification. In other words, depending on the way the system was modelled, making small changes to the specification such as adding new features to the software may require significant rewrites of the proofs \cite{fv-drift}.
\\

\subsection{Interactive Theorem Provers}
Interactive theorem provers (also known as Proof Assistants) are tools that assists users in constructing formal proofs in a step-by-step manner, verifying their proof against a set of logical rules, often backed by some form of higher-order logic system. This allows users to write proofs that can be computer-verified to be correct, improving confidence in their formally verified system.
\\

Interactive theorem provers may vary in their underlying logical foundations. For example, Coq and Lean use the Calculus of Inductive Constructions \cite{coq, lean4}, Agda relies on an extension of Martin-Löf type theory \cite{agda}, while Idris is based on Quantitative Type Theory \cite{idris2}. Despite being a relatively recent development, interactive theorem provers have been used to formally verify notable mathematical results, such as a proof for the Four-Color Theorem \cite{4color} in Coq, a proof that $BB(5) = 47176870$ in Coq \cite{busybeaver}, and a proof of the Polynomial Freiman-Rusza Conjecture in Lean 4 \cite{pfr-lean}.
\\

\subsection{Formally Verifying Byzantine Fault-Tolerant Consensus Protocols}

Given the importance of BFT consensus protocols, it would be sensible to formally verify such protocols to ensure that the requisite properties hold under reasonable assumptions \cite{dsbugs}. Bythos is one such work that provides a framework for users to formally verify their own custom protocols in an \textit{modular} and \textit{incremental} manner \cite{bythos}. Bythos, in the authors' own words, functions in a "pay as you go" manner, where proofs of stronger properties can be built up from proofs of weaker properties. 
\\

To illustrate how this is possible, we cite a motivating example for Bythos, which is that of iterated Provable Broadcast (PB). Provable Broadcast is a simple BFT protocol that guarantees certain properties, such as \textit{Termination}, \textit{Uniqueness}, \textit{External Validity}, and \textit{Weak Availability}. However, by \textbf{iterating} Provable Broadcast (i.e. running it multiple times sequentially before taking the final output), we can guarantee increasingly stronger properties on the resultant protocol. For example, iterating PB twice yields the Locked Broadcast protocol, which guarantees \textit{Unique-Lock-Availability}; iterating thrice yields Keyed Broadcast, which guarantees \textit{Unique-Key-Availability}; and iterating four times yields Robust Keyed Broadcast, which guarantees \textit{Robust-Delivery} \cite{pb}. Note that these iterated variants of Provable Broadcast are also used as building blocks for other modern BFT protocols, such as Tendermint, HotStuff, and VABA. By building up our proofs for these stronger guarantees over each iteration, Bythos allows users to see efficiency gains as we avoid having to begin verification from scratch with each new derived protocol.
\\

\subsection{Refinement and Extraction}
In our earlier discussion on formal verification, we describe how we start with an existing implementation of a system and proceed to construct a mathematical model based on that system. We also observed that our model should be a faithful representation of the given implementation, otherwise our properties may no longer be guaranteed on our concrete implementation.
\\

However, working in the other direction is also possible. Using a technique called Program Refinement, we can start from an abstraction (or mathematical model) of the specification, and progressively reduce the level of abstraction by adding concrete elements (such as implementation details like data structures or optimizations), while proving that such details do not violate the invariants / properties proven on the original model \cite{refinement}. This ultimately allows us to create a concrete implementation that is guaranteed to match our mathematical model and thus fulfills the required specification.
\\

Note that developing the refined program may not necessarily be an automated process. In some cases, we are able to directly \textit{extract} a program from the original model. One key example of such Program Extraction is with the Coq proof assistant, where any structures modelled in the Coq language can be automatically transformed into equivalent programs in more popular languages like OCaml, Haskell or Scheme \cite{extraction}.
\\

One limitation of extraction is that the generated code may be difficult to read and understand by human developers, which may make it challenging to implement low-level performance optimizations. For example, Coq-extracted OCaml code may include liberal use of \code{Obj.magic}, a type-casting function which is normally avoided among OCaml developers for safety reasons. However, such cases are permissible here, since the program was verified in Coq, which has a stronger type system than OCaml. Nevertheless, the rationale for such "unsafe" casts may not be clear to the programmer just from reading the extracted code.

~\\

\section{Rust} 
When developing a library as a front-end over protocols implemented in Lean, there are a wide variety of mainstream programming languages that we can choose as the target language for our library. Ultimately, we chose to develop our library in the Rust programming language.
\\

Rust is a systems programming language created at Mozilla. The language's main advantages are in its performance, type safety, and concurrency, which is achieved in no small part due to its memory model \cite{rustlang}. Rust is not garbage-collected, instead, it uses an ownership model to keep track of the lifetimes of variables, allowing the runtime to free the memory used by variables as soon as their lifetime ends. This ownership system is also responsible for preventing data races by allowing a variable to either be referenced by multiple readers or mutated by a single writer at any given point in time. Most importantly, this ownership system is embedded in its type system, which is in turn enforced by Rust's compiler. As a result, these checks are done at compile-time and do not contribute to runtime bloat, allowing programs written in Rust to be both safe and fast. 
\\

In addition, Rust also has a well-established package ecosystem, with many additional packages (known as crates) available on the repository crates.io. Notably, Rust has several battle-tested networking and asynchronous programming crates available \cite{tokio, libp2p}. As a result, Rust has gained increased adoption particularly within the blockchain community,  which is a strong pull factor towards our adoption of Rust as the target language of choice for our library.
\\

~\\